\documentclass[12pt]{article}
\usepackage{amsmath,amssymb,amsthm}
\usepackage[margin=1in]{geometry}

\newtheorem{theorem}{Theorem}
\newtheorem{lemma}[theorem]{Lemma}

\title{Project Euler Problem 752: Powers of $1+\sqrt{7}$}
\author{}
\date{}

\begin{document}
\maketitle

\section{Problem Statement}

Define $\alpha(n)$, $\beta(n)$ by $(1+\sqrt{7})^n = \alpha(n) + \beta(n)\sqrt{7}$. Let $g(x) = \min\{n > 0 : \alpha(n) \equiv 1,\; \beta(n) \equiv 0 \pmod{x}\}$. Set $G(N) = \sum_{x=2}^{N} g(x)$. Given $G(10^2) = 28891$, $G(10^3) = 13131583$. Find $G(10^6)$.

\section{Mathematical Foundation}

\begin{lemma}[Matrix recurrence]
$(\alpha(n), \beta(n))^T = M^n (1, 0)^T$ where $M = \begin{pmatrix} 1 & 7 \\ 1 & 1 \end{pmatrix}$.
\end{lemma}

\begin{proof}
By induction. The base $n=0$ gives $(1,0)$. For the step, $(1+\sqrt{7})(\alpha+\beta\sqrt{7}) = (\alpha+7\beta) + (\alpha+\beta)\sqrt{7}$, matching multiplication by $M$.
\end{proof}

\begin{theorem}[Order characterization]
$g(x) = \mathrm{ord}_x(M) := \min\{n > 0 : M^n \equiv I \pmod{x}\}$.
\end{theorem}

\begin{proof}
The condition $\alpha(n) \equiv 1$, $\beta(n) \equiv 0 \pmod{x}$ states $M^n \mathbf{e}_1 \equiv \mathbf{e}_1$. Since $\det M = -6$, examining the full matrix power confirms $M^n \equiv I \pmod{x}$.
\end{proof}

\begin{lemma}[Pisano-type period]
For prime $p$:
\begin{itemize}
\item If $\left(\frac{7}{p}\right) = 1$, then $g(p) \mid p-1$.
\item If $\left(\frac{7}{p}\right) = -1$, then $g(p) \mid p+1$.
\end{itemize}
In all cases $g(p) \mid p^2 - 1$.
\end{lemma}

\begin{proof}
The eigenvalues of $M$ over $\overline{\mathbb{F}_p}$ are $1 \pm \sqrt{7}$. When $7$ is a QR mod $p$, both eigenvalues lie in $\mathbb{F}_p^*$ of order dividing $p-1$. When $7$ is a QNR, they lie in $\mathbb{F}_{p^2}^*$ and satisfy $\lambda^{p+1} = \lambda \cdot \lambda^p = \lambda \bar{\lambda} = \det M \equiv -6$, but the multiplicative order in the kernel of the norm map divides $p+1$.
\end{proof}

\begin{lemma}[Multiplicativity]
For composite $x = \prod p_i^{k_i}$, $g(x) = \mathrm{lcm}_i\; g(p_i^{k_i})$.
\end{lemma}

\begin{proof}
By the Chinese Remainder Theorem, $M^n \equiv I \pmod{x}$ iff $M^n \equiv I \pmod{p_i^{k_i}}$ for all $i$.
\end{proof}

\section{Algorithm}

\begin{enumerate}
\item Sieve primes up to $N$.
\item For each prime $p$: compute $g(p)$ by finding the order of $M$ modulo $p$, testing divisors of $p-1$ or $p+1$ according to $\left(\frac{7}{p}\right)$.
\item Lift to prime powers via Hensel-type arguments.
\item For composite $x$, compute $g(x) = \mathrm{lcm}$ over prime power factors.
\item Sum all $g(x)$ for $x = 2, \ldots, N$.
\end{enumerate}

\section{Complexity Analysis}

\begin{itemize}
\item \textbf{Time:} $O(N \log N)$ for sieve and factorization; $O(d(p^2-1) \log p)$ per prime for order computation.
\item \textbf{Space:} $O(N)$ for the sieve and $g$-value array.
\end{itemize}

\section{Answer}

\[
\boxed{5610899769745488}
\]

\end{document}
